% ============================================================
\chapter{Th\'eorie des Groupes pour le Deep Learning}
\label{ch:groupes}
% ============================================================

Les sym\'etries sont partout~: un visage est (approximativement) sym\'etrique
par r\'eflexion, un cristal l'est par rotation, les lois de la physique sont
invariantes par translation dans l'espace et le temps. La th\'eorie des
groupes, n\'ee des travaux de Galois et d'Abel au XIX\ieme{} si\`ecle, est le
langage math\'ematique des sym\'etries. Et en deep learning, exploiter les
sym\'etries d'un probl\`eme --- les int\'egrer comme \emph{biais inductif} dans
l'architecture du r\'eseau --- est devenu l'une des id\'ees les plus
puissantes de la derni\`ere d\'ecennie.

% ------------------------------------------------------------
\section{Introduction \`a la th\'eorie des groupes}
% ------------------------------------------------------------

La th\'eorie des groupes fournit le langage math\'ematique pour d\'ecrire les
sym\'etries. Ce chapitre introduit les concepts fondamentaux n\'ecessaires \`a
la compr\'ehension de l'apprentissage g\'eom\'etrique profond.

\begin{definition}[Groupe]
Un \textbf{groupe} est un ensemble $G$ muni d'une op\'eration binaire
$\cdot: G \times G \to G$ satisfaisant~:
\begin{enumerate}
    \item \textbf{Associativit\'e}~: $(g \cdot h) \cdot k = g \cdot (h \cdot k)$
          pour tous $g, h, k \in G$.
    \item \textbf{\'El\'ement neutre}~: il existe $e \in G$ tel que
          $e \cdot g = g \cdot e = g$ pour tout $g \in G$.
    \item \textbf{Inverse}~: pour tout $g \in G$, il existe $g^{-1} \in G$ tel que
          $g \cdot g^{-1} = g^{-1} \cdot g = e$.
\end{enumerate}
\end{definition}

\begin{exemple}[Groupes fondamentaux en GDL]
\begin{enumerate}
    \item \textbf{Groupe sym\'etrique} $S_n$~: permutations de $n$ \'el\'ements.
          $\abs{S_n} = n!$.
    \item \textbf{Groupe cyclique} $\Z_n = \Z/n\Z$~: rotations discr\`etes.
    \item \textbf{Groupe orthogonal} $\mathrm{O}(n)$~: matrices $\mathbf{Q} \in \R^{n \times n}$
          telles que $\mathbf{Q}^\top\mathbf{Q} = \mathbf{I}$.
    \item \textbf{Groupe sp\'ecial orthogonal} $\mathrm{SO}(n)$~: rotations,
          $\det(\mathbf{Q}) = +1$.
    \item \textbf{Groupe euclidien} $\mathrm{E}(n) = \mathrm{O}(n) \ltimes \R^n$~:
          rotations, r\'eflexions et translations.
    \item \textbf{Groupe sp\'ecial euclidien} $\mathrm{SE}(n) = \mathrm{SO}(n) \ltimes \R^n$~:
          rotations et translations (mouvements rigides).
\end{enumerate}
\end{exemple}

\subsection{Sous-groupes et homomorphismes}

\begin{definition}[Sous-groupe]
Un sous-ensemble $H \subseteq G$ est un \textbf{sous-groupe} de $G$ (not\'e
$H \leq G$) si $H$ est lui-m\^eme un groupe pour l'op\'eration de $G$.
\end{definition}

\begin{definition}[Homomorphisme de groupes]
Un \textbf{homomorphisme} de groupes est une application $\varphi: G \to H$
telle que $\varphi(g_1 \cdot g_2) = \varphi(g_1) \cdot \varphi(g_2)$ pour
tous $g_1, g_2 \in G$. Si $\varphi$ est bijectif, c'est un \textbf{isomorphisme}.
\end{definition}

\begin{proposition}[Propri\'et\'es des homomorphismes]
Soit $\varphi: G \to H$ un homomorphisme. Alors~:
\begin{enumerate}
    \item $\varphi(e_G) = e_H$.
    \item $\varphi(g^{-1}) = \varphi(g)^{-1}$ pour tout $g \in G$.
    \item $\ker(\varphi) = \{g \in G : \varphi(g) = e_H\}$ est un sous-groupe
          normal de $G$.
\end{enumerate}
\end{proposition}

\section{Actions de groupe}

\begin{definition}[Action de groupe]
Une \textbf{action} (gauche) d'un groupe $G$ sur un ensemble $X$ est une
application $\cdot: G \times X \to X$ telle que~:
\begin{enumerate}
    \item $e \cdot x = x$ pour tout $x \in X$.
    \item $(g \cdot h) \cdot x = g \cdot (h \cdot x)$ pour tous $g, h \in G$
          et $x \in X$.
\end{enumerate}
\end{definition}

\begin{exemple}[Actions en deep learning]
\begin{enumerate}
    \item $S_n$ agit sur $\R^n$ par permutation des coordonn\'ees~:
          $\pi \cdot (x_1, \ldots, x_n) = (x_{\pi^{-1}(1)}, \ldots, x_{\pi^{-1}(n)})$.
    \item $\mathrm{SO}(3)$ agit sur $\R^3$ par rotation~: $R \cdot \mathbf{x} = R\mathbf{x}$.
    \item $\Z^2$ agit sur les images par translation~: $\mathbf{t} \cdot f(\mathbf{x}) = f(\mathbf{x} - \mathbf{t})$.
\end{enumerate}
\end{exemple}

\begin{definition}[Orbite et stabilisateur]
Soit $G$ agissant sur $X$. Pour $x \in X$~:
\begin{itemize}
    \item L'\textbf{orbite} de $x$ est $G \cdot x = \{g \cdot x : g \in G\}$.
    \item Le \textbf{stabilisateur} de $x$ est $G_x = \{g \in G : g \cdot x = x\}$.
\end{itemize}
\end{definition}

\begin{theoreme}[Th\'eor\`eme orbite-stabilisateur]
Pour un groupe fini $G$ agissant sur $X$ et $x \in X$~:
\[
    \abs{G} = \abs{G \cdot x} \cdot \abs{G_x}
\]
\end{theoreme}

\section{Repr\'esentations de groupes}

\begin{definition}[Repr\'esentation]
Une \textbf{repr\'esentation} d'un groupe $G$ sur un espace vectoriel $V$ est
un homomorphisme $\rho: G \to \mathrm{GL}(V)$, c'est-\`a-dire une application
qui associe \`a chaque $g \in G$ une transformation lin\'eaire inversible
$\rho(g): V \to V$ telle que $\rho(g_1 g_2) = \rho(g_1)\rho(g_2)$.
\end{definition}

\begin{exemple}[Repr\'esentations courantes]
\begin{enumerate}
    \item \textbf{Repr\'esentation triviale}~: $\rho(g) = \mathrm{Id}$ pour tout $g$.
    \item \textbf{Repr\'esentation par permutation}~: $\rho(\pi) = \mathbf{P}_\pi$
          (matrice de permutation).
    \item \textbf{Repr\'esentation standard de $\mathrm{SO}(3)$}~: $\rho(R) = R$
          (la matrice de rotation elle-m\^eme).
    \item \textbf{Repr\'esentation r\'eguli\`ere}~: $[\rho(g)f](h) = f(g^{-1}h)$.
\end{enumerate}
\end{exemple}

\subsection{Repr\'esentations irr\'eductibles}

\begin{definition}[Repr\'esentation irr\'eductible]
Une repr\'esentation $\rho: G \to \mathrm{GL}(V)$ est \textbf{irr\'eductible}
si les seuls sous-espaces de $V$ invariants sous l'action de $G$ sont $\{0\}$
et $V$ lui-m\^eme.
\end{definition}

\begin{theoreme}[D\'ecomposition de Maschke]
Toute repr\'esentation de dimension finie d'un groupe fini (sur $\R$ ou $\C$)
peut \^etre d\'ecompos\'ee en somme directe de repr\'esentations irr\'eductibles~:
\[
    V = V_1 \oplus V_2 \oplus \cdots \oplus V_k
\]
o\`u chaque $V_i$ est un sous-espace irr\'eductible.
\end{theoreme}

\begin{formules}[title=\textbf{Repr\'esentations irr\'eductibles de $\mathrm{SO}(3)$}]
Les repr\'esentations irr\'eductibles de $\mathrm{SO}(3)$ sont index\'ees par
$\ell \in \{0, 1, 2, \ldots\}$ et ont dimension $2\ell + 1$~:
\begin{itemize}
    \item $\ell = 0$~: scalaires (dimension 1) --- repr\'esentation triviale.
    \item $\ell = 1$~: vecteurs (dimension 3) --- repr\'esentation standard.
    \item $\ell = 2$~: tenseurs sym\'etriques sans trace (dimension 5).
\end{itemize}
Les matrices de repr\'esentation sont les \textbf{matrices de Wigner}
$D^\ell_{mm'}(R)$, reli\'ees aux harmoniques sph\'eriques~:
\[
    Y_\ell^m(R^{-1}\hat{\mathbf{r}}) = \sum_{m'=-\ell}^{\ell}
    D^\ell_{mm'}(R)\, Y_\ell^{m'}(\hat{\mathbf{r}})
\]
\end{formules}

\section{\'Equivariance et invariance}

\begin{definition}[Application \'equivariante]
Soient $\rho_1: G \to \mathrm{GL}(V_1)$ et $\rho_2: G \to \mathrm{GL}(V_2)$
deux repr\'esentations. Une application lin\'eaire $T: V_1 \to V_2$ est
\textbf{\'equivariante} (ou \textbf{entrelacement}) si~:
\[
    T \circ \rho_1(g) = \rho_2(g) \circ T, \quad \forall g \in G.
\]
Lorsque $\rho_2$ est la repr\'esentation triviale, $T$ est \textbf{invariante}.
\end{definition}

\begin{theoreme}[Lemme de Schur]
Soient $\rho_1$ et $\rho_2$ deux repr\'esentations irr\'eductibles de $G$ et
$T: V_1 \to V_2$ un entrelacement~:
\begin{enumerate}
    \item Si $\rho_1 \not\cong \rho_2$, alors $T = 0$.
    \item Si $\rho_1 \cong \rho_2$ (sur $\C$), alors $T = \lambda\, \mathrm{Id}$
          pour un certain $\lambda \in \C$.
\end{enumerate}
\end{theoreme}

\begin{intuition}[title=\textbf{Importance du lemme de Schur}]
Le lemme de Schur est fondamental car il \textbf{contraint fortement} la forme
des couches \'equivariantes dans un r\'eseau de neurones. Si l'on impose
l'\'equivariance, l'espace des applications lin\'eaires autoris\'ees est
drastiquement r\'eduit, ce qui justifie le partage des poids.
\end{intuition}

\section{Produit tensoriel de repr\'esentations}

\begin{definition}[Produit tensoriel]
Le \textbf{produit tensoriel} de deux repr\'esentations $\rho_1: G \to \mathrm{GL}(V_1)$
et $\rho_2: G \to \mathrm{GL}(V_2)$ est la repr\'esentation
$\rho_1 \otimes \rho_2: G \to \mathrm{GL}(V_1 \otimes V_2)$ d\'efinie par~:
\[
    (\rho_1 \otimes \rho_2)(g)(v_1 \otimes v_2) = \rho_1(g)v_1 \otimes \rho_2(g)v_2.
\]
\end{definition}

\begin{theoreme}[D\'ecomposition de Clebsch--Gordan pour $\mathrm{SO}(3)$]
Le produit tensoriel de deux repr\'esentations irr\'eductibles de $\mathrm{SO}(3)$
se d\'ecompose selon~:
\[
    D^{\ell_1} \otimes D^{\ell_2} = \bigoplus_{\ell=\abs{\ell_1-\ell_2}}^{\ell_1+\ell_2} D^{\ell}
\]
Les coefficients de changement de base sont les \textbf{coefficients de
Clebsch--Gordan} $C^{\ell m}_{\ell_1 m_1, \ell_2 m_2}$.
\end{theoreme}

\begin{codeblock}[title=\textbf{Produit tensoriel avec e3nn}]
\begin{minted}{python}
import torch
from e3nn import o3

# Representations irreductibles de SO(3)
irreps_1 = o3.Irreps("1x0e + 1x1o")  # scalaire + vecteur
irreps_2 = o3.Irreps("1x1o")          # vecteur

# Produit tensoriel
tp = o3.FullTensorProduct(irreps_1, irreps_2)
print(f"Entree 1: {irreps_1}")
print(f"Entree 2: {irreps_2}")
print(f"Sortie: {tp.irreps_out}")
# 0e x 1o -> 1o, 1o x 1o -> 0e + 1e + 2e

x1 = irreps_1.randn(5, -1)  # batch de 5
x2 = irreps_2.randn(5, -1)
out = tp(x1, x2)
print(f"Shape sortie: {out.shape}")
\end{minted}
\end{codeblock}

\section{Groupes de Lie}

\begin{definition}[Groupe de Lie]
Un \textbf{groupe de Lie} est un groupe $G$ qui est aussi une vari\'et\'e
diff\'erentiable, de sorte que les op\'erations de multiplication
$(g, h) \mapsto gh$ et d'inversion $g \mapsto g^{-1}$ sont lisses.
\end{definition}

\begin{exemple}[Groupes de Lie en GDL]
\begin{center}
\begin{tabular}{llll}
\toprule
\textbf{Groupe} & \textbf{Dimension} & \textbf{Compact~?} & \textbf{Usage} \\
\midrule
$\mathrm{SO}(2)$ & 1 & Oui & Rotations 2D \\
$\mathrm{SO}(3)$ & 3 & Oui & Rotations 3D \\
$\mathrm{SE}(3)$ & 6 & Non & Mouvements rigides \\
$\mathrm{O}(3)$ & 3 & Oui & Rotations + r\'eflexions \\
$\mathrm{GL}(n, \R)$ & $n^2$ & Non & Transformations lin\'eaires \\
\bottomrule
\end{tabular}
\end{center}
\end{exemple}

\subsection{Alg\`ebre de Lie}

\begin{definition}[Alg\`ebre de Lie]
L'\textbf{alg\`ebre de Lie} $\mathfrak{g}$ d'un groupe de Lie $G$ est l'espace
tangent \`a $G$ en l'\'el\'ement neutre $e$, muni du \textbf{crochet de Lie}
$[\cdot, \cdot]: \mathfrak{g} \times \mathfrak{g} \to \mathfrak{g}$.
\end{definition}

\begin{formules}[title=\textbf{Application exponentielle}]
L'\textbf{application exponentielle} $\exp: \mathfrak{g} \to G$ relie
l'alg\`ebre de Lie au groupe~:
\[
    \exp(X) = \sum_{k=0}^{\infty} \frac{X^k}{k!}
\]
Pour $\mathrm{SO}(3)$, l'alg\`ebre de Lie $\mathfrak{so}(3)$ est l'espace des
matrices antisym\'etriques $3 \times 3$. La formule de Rodrigues donne~:
\[
    \exp(\theta\, [\hat{\mathbf{n}}]_\times) = \mathbf{I} + \sin\theta\,
    [\hat{\mathbf{n}}]_\times + (1 - \cos\theta)\, [\hat{\mathbf{n}}]_\times^2
\]
o\`u $[\hat{\mathbf{n}}]_\times$ est la matrice antisym\'etrique associ\'ee au
vecteur unitaire $\hat{\mathbf{n}}$.
\end{formules}

\section{Convolution sur les groupes}

\begin{definition}[Convolution de groupe]
Soit $G$ un groupe localement compact muni d'une mesure de Haar $\mu$. La
\textbf{convolution de groupe} de deux fonctions $f, g: G \to \R$ est~:
\[
    (f * g)(x) = \int_G f(y)\, g(y^{-1}x)\, d\mu(y).
\]
\end{definition}

\begin{proposition}[\'Equivariance de la convolution]
La convolution de groupe est \'equivariante par translation \`a gauche~: si
$L_a f(x) = f(a^{-1}x)$, alors~:
\[
    L_a(f * g) = (L_a f) * g.
\]
C'est cette propri\'et\'e qui justifie l'utilisation de la convolution comme
couche fondamentale dans les r\'eseaux \'equivariants.
\end{proposition}

\begin{theoreme}[CNN comme convolution sur $\Z^2$]
Un r\'eseau convolutif classique r\'ealise une convolution de groupe sur
$G = \Z^2$ (le groupe de translation discret). La premi\`ere couche effectue
une corr\'elation crois\'ee~:
\[
    [f \star \psi](x) = \sum_{y \in \Z^2} f(y)\, \psi(y - x)
    = \sum_{y \in \Z^2} f(x + y)\, \psi(y)
\]
qui est \'equivariante par translation de $f$.
\end{theoreme}

\section{Group Equivariant CNNs (G-CNNs)}

Cohen et Welling (2016) ont g\'en\'eralis\'e les CNN en rempla\c{c}ant le
groupe de translation $\Z^2$ par un groupe plus grand $G$.

\begin{definition}[G-CNN]
Un \textbf{G-CNN} est un r\'eseau dont les couches r\'ealisent des convolutions
sur un groupe $G$. La premi\`ere couche \og souleve \fg (\textit{lifts}) le
signal de $\Z^2$ vers $G$~:
\[
    [\text{lift}(f)](g) = \sum_{x \in \Z^2} f(x)\, \psi(g^{-1} \cdot x)
\]
Les couches suivantes op\`erent enti\`erement sur $G$~:
\[
    [f * \psi](g) = \sum_{h \in G} f(h)\, \psi(h^{-1}g)
\]
\end{definition}

\begin{codeblock}[title=\textbf{G-CNN pour le groupe $p4$ (rotations 90\textdegree)}]
\begin{minted}{python}
import torch
import torch.nn as nn

class P4Conv(nn.Module):
    """Convolution equivariante par le groupe p4 (rotations 0/90/180/270)."""
    def __init__(self, in_channels, out_channels, kernel_size=3):
        super().__init__()
        self.weight = nn.Parameter(
            torch.randn(out_channels, in_channels, kernel_size, kernel_size)
        )
        nn.init.kaiming_normal_(self.weight)

    def _rotate_kernel(self, w, k):
        """Rotation du noyau de k*90 degres."""
        return torch.rot90(w, k, dims=[-2, -1])

    def forward(self, x):
        # x: (batch, C_in, H, W) ou (batch, C_in, 4, H, W)
        outputs = []
        for k in range(4):  # 4 rotations
            w_rot = self._rotate_kernel(self.weight, k)
            out = torch.nn.functional.conv2d(
                x if x.dim() == 4 else x.sum(dim=2),
                w_rot, padding=1
            )
            outputs.append(out)
        # Stack sur la dimension du groupe
        return torch.stack(outputs, dim=2)  # (batch, C_out, 4, H, W)

# Test
conv = P4Conv(3, 16)
x = torch.randn(2, 3, 32, 32)
out = conv(x)
print(f"Sortie: {out.shape}")  # (2, 16, 4, 32, 32)
\end{minted}
\end{codeblock}

\section{Invariants et \'equivariants lin\'eaires}

\begin{theoreme}[Caract\'erisation des applications lin\'eaires \'equivariantes]
L'espace des applications lin\'eaires $T: V_1 \to V_2$ qui sont \'equivariantes
par rapport aux repr\'esentations $\rho_1$ et $\rho_2$ est~:
\[
    \mathrm{Hom}_G(V_1, V_2) = \{T \in \mathrm{Hom}(V_1, V_2) :
    T\rho_1(g) = \rho_2(g)T,\; \forall g \in G\}
\]
Sa dimension est donn\'ee par~:
\[
    \dim \mathrm{Hom}_G(V_1, V_2) = \inner{\chi_1}{\chi_2}_G
    = \frac{1}{\abs{G}} \sum_{g \in G} \overline{\chi_1(g)}\, \chi_2(g)
\]
o\`u $\chi_i(g) = \tr(\rho_i(g))$ est le \textbf{caract\`ere} de $\rho_i$.
\end{theoreme}

\begin{corollaire}[Nombre de param\`etres libres]
Pour une couche \'equivariante entre deux repr\'esentations d\'ecompos\'ees en
irr\'eductibles $V_1 = \bigoplus_i n_i V_i^{\text{irr}}$ et
$V_2 = \bigoplus_j m_j V_j^{\text{irr}}$, le nombre de param\`etres libres est~:
\[
    \sum_k n_k \cdot m_k
\]
o\`u la somme porte sur les types irr\'eductibles communs.
\end{corollaire}

\section*{Exercices}

\begin{exercice}[Groupes de sym\'etrie de mol\'ecules]
\begin{enumerate}
    \item D\'eterminer le groupe de sym\'etrie de la mol\'ecule d'eau $\mathrm{H_2O}$
          (groupe $C_{2v}$).
    \item \'Enum\'erer toutes les repr\'esentations irr\'eductibles de $C_{2v}$.
    \item Expliquer comment ces sym\'etries contraignent un mod\`ele de pr\'ediction
          d'\'energie mol\'eculaire.
\end{enumerate}
\end{exercice}

\begin{exercice}[Repr\'esentations de $S_3$]
\begin{enumerate}
    \item \'Enum\'erer les six \'el\'ements du groupe sym\'etrique $S_3$.
    \item Trouver toutes les repr\'esentations irr\'eductibles de $S_3$ (il y en a trois).
    \item V\'erifier la relation $\sum_i (\dim V_i)^2 = \abs{G}$.
    \item D\'ecomposer la repr\'esentation par permutation standard $\R^3$ en
          irr\'eductibles.
\end{enumerate}
\end{exercice}

\begin{exercice}[Formule de Rodrigues]
\begin{enumerate}
    \item V\'erifier la formule de Rodrigues pour $\theta = \pi/2$ et
          $\hat{\mathbf{n}} = (0, 0, 1)$.
    \item Montrer que $\exp(\theta [\hat{\mathbf{n}}]_\times) \in \mathrm{SO}(3)$.
    \item Impl\'ementer l'application exponentielle $\mathfrak{so}(3) \to \mathrm{SO}(3)$
          en PyTorch et v\'erifier l'orthogonalit\'e du r\'esultat.
\end{enumerate}
\end{exercice}

\begin{exercice}[Convolution de groupe discr\`ete]
Impl\'ementer une convolution de groupe sur le groupe di\'edral $D_4$
(sym\'etries du carr\'e).
\begin{enumerate}
    \item \'Enum\'erer les 8 \'el\'ements de $D_4$.
    \item Impl\'ementer la table de multiplication de $D_4$.
    \item Impl\'ementer la convolution de groupe et v\'erifier son \'equivariance.
\end{enumerate}
\end{exercice}
